% chapter2-beamer.tex — a standalone slide deck for Chapter 2.
% See chapter1/beamer/chapter1-beamer.tex for the full explanation of the
% `slim` option and the adapt-don't-\input approach.  Build with
% `make chapter2-beamer`.

\documentclass{beamer}

\usetheme{Madrid}
\setbeamertemplate{navigation symbols}{}
\setbeamertemplate{theorems}[numbered]

\usepackage[slim]{philogic}

\usepackage[backend=biber]{biblatex}
\addbibresource{../../refs.bib}

% beamer predefines theorem, corollary, definition, example, and proof — but not
% proposition.  In slim mode philogic doesn't add it either (the host class is
% in charge of theorem environments), so define the one we need here:
\newtheorem{proposition}{Proposition}

\title{Frame Conditions \& Neighbourhood Semantics}
\subtitle{Chapter 2}
\author{Your Name}
\institute{Your Department \\ University of California, Davis}
\date{}

\begin{document}

\frame{\titlepage}

\begin{frame}{Frame conditions}
  Properties of the accessibility relation correspond to modal axioms.
  \begin{proposition}
    A frame \FR{} validates \tax{} ($\nec A \to A$) if and only if $R$ is
    reflexive.
  \end{proposition}
  Likewise, transitivity corresponds to $\nec A \to \nec\nec A$ and symmetry to
  $A \to \nec\poss A$ \parencite[ch.~3]{blackburn2001}.
\end{frame}

\begin{frame}{Neighbourhood semantics}
  Generalise relational models by replacing $R$ with a neighbourhood function
  \parencite{scott1970, pacuit2017}.
  \begin{definition}
    A \emph{neighbourhood model} $\MN = \oset{W, \fn, V}$ has
    $\fn : W \to \power{\power{W}}$, assigning each world a collection of
    neighbourhoods.  The modal clause:
    \[
      \MN,w \trues \nec A
      \quad\text{iff}\quad
      \truthset{A}^{\MN} \in \fn(w).
    \]
  \end{definition}
  \begin{block}{Why bother?}
    Every Kripke frame induces a neighbourhood frame, but not conversely:
    neighbourhood frames model non-normal logics lacking \kax{} or \nrule{}
    \parencite{chellas1980}.
  \end{block}
\end{frame}

\begin{frame}[allowframebreaks]{References}
  \printbibliography
\end{frame}

\end{document}
